\documentclass[10pt,a4paper]{ctexart}
\usepackage[margin=2cm]{geometry}
\usepackage{graphicx}
\usepackage{hyperref}
\usepackage{bookmark}
\usepackage{amsmath}    % 提供方程环境
\usepackage{physics}    % 简化向量、微分算子语法
\newcommand{\myspace}[1]{\par\vspace{#1\baselineskip}} %自定义空一行
\usepackage{fancyhdr}
\pagestyle{fancy}
\fancyhf{}  % 清除所有页眉页脚
\fancyfoot[C]{\thepage}  % 在页脚中央设置页码
\renewcommand{\headrulewidth}{0pt}  % 去掉页眉横线
\renewcommand{\footrulewidth}{0pt}  % 去掉页脚横线

\usepackage{hyperref}
\hypersetup{
    colorlinks=true,
    linkcolor=blue,
    filecolor=magenta,      
    urlcolor=cyan,
    pdftitle={Overleaf Example},
    pdfpagemode=FullScreen,
    }

\usepackage{titlesec}

% 重新定义section格式
\titleformat{\section}
  {\normalfont\Large\bfseries}
  {\chinese{section}、}
  {0pt}
  {}

% 给出常用字号的自定义指令
\newcommand{\chuhao}{\fontsize{42pt}{48pt}\selectfont}
\newcommand{\xiaochu}{\fontsize{36pt}{42pt}\selectfont}
\newcommand{\xiaoer}{\fontsize{18pt}{21pt}\selectfont}
\newcommand{\sanhao}{\fontsize{16pt}{20pt}\selectfont}
\newcommand{\xiaosan}{\fontsize{15pt}{18pt}\selectfont}
\newcommand{\sihao}{\fontsize{14pt}{18pt}\selectfont}
\newcommand{\xiaosi}{\fontsize{12pt}{16pt}\selectfont}
\newcommand{\wuhao}{\fontsize{11pt}{12pt}\selectfont}

% 定义中文数字转换
\titleformat{\section}
  {\normalfont\Large\bfseries}
  {\zhnum{section}、} 
  {0pt}
  {}

% 标题设置
\title{\xiaochu\textbf{物理实验预习报告}}
\author{}
\date{}

\begin{document}

\vspace*{36pt} % 顶部空白

% 居中填写区域
\begin{center}


    \chuhao\textbf{物理实验预习报告}
    % 预留空白区域
    \vspace{13pt}
    \vspace{13pt}
    \vspace{13pt}
    \vspace{13pt}
    
    \sanhao\textbf{实验名称：\underline{\makebox[10cm]{ 光电效应测定普朗克常量}}} \\[10pt]
    \sanhao\textbf{指导教师：\underline{\makebox[10cm]{ 史明 }}} \\[30pt]
    
    % 预留空白区域
    \myspace{7}
    
   \sihao\textbf{班级：\underline{\makebox[6cm]{}}} \\[10pt]
    \sihao\textbf{姓名：\underline{\makebox[6cm]{}}} \\[10pt]
    \sihao\textbf{学号：\underline{\makebox[6cm]{}}} \\[30pt]
    
    \sihao\textbf{实验日期: \underline{\makebox[0.8cm]{2025}}年 \underline{\makebox[0.4cm]{10}}月\underline{\makebox[0.4cm]{15}}日  星期\underline{\makebox[0.6cm]{三}}上午}
    
    \vspace{18pt}
    \sihao 浙江大学物理实验教学中心
    
\end{center}

\newpage

\section{实验综述}

\xiaosi （自述实验现象、实验原理和实验方法，不超过300字，5分）

光电效应是光照射金属表面时电子逸出的现象。
实验现象表明：存在截止频率$\nu_0$，低于此频率无光电流；光电子的最大初动能与光强无关。
依据爱因斯坦光电效应方程$h\nu = \frac{1}{2}mv_m^2 + W$ 和$eU_0=\frac{1}{2}mv^2_m$ 有$U_0=\frac{h}{e} \nu-\frac{W}{e}$，其中$W$为逸出功。
实验通过测量不同频率$\nu$单色光照射下的遏止电压$U_0$，验证$U_0-\nu$的线性关系：$U_0 = \frac{h}{e}\nu - \frac{W}{e}$。
使用汞灯光源配合滤光片获得单色光，调节反向电压使光电流为零，此时电压即为$U_0$。
绘制$U_0-\nu$曲线，其斜率$k = h/e$，可计算普朗克常量$h = ke$。

此外，通过测量光电管电压和电流可绘制光电管的伏安特性曲线，测量光源距离、光阑孔径和饱和光电流可以验证饱和光电流与光强的正比关系。
\myspace{2}

\section{实验重点}

\xiaosi（简述本实验的学习重点，不超过100字，3分）

1. 精确测量不同频率光对应的遏止电压$U_0$ 
 
2. 准确绘制$U_0-\nu$图像并计算斜率，通过斜率$k$计算普朗克常量$h = ke$ 

3. 掌握消除暗电流和本底电流影响的方法
\myspace{2}

\section{实验难点}

\xiaosi（简述本实验的实现难点，不超过100字，2分）

1. 遏止电压$U_0$的准确判定：实际光电流为零点难以确定 

2. 暗电流和反向电流的干扰，需使用方法减少其影响和干扰 

3. 杂散光影响测量精度 

4. 光电管疲劳效应引起的灵敏度变化


\vspace{3cm}

\sihao\textbf{注意事项：}
\xiaosi
\begin{enumerate}
    \item 用PDF格式上传“预习报告”，文件名：学生姓名+学号+实验名称+周次。
    \item “预习报告”必须递交在“学在浙大”的本课程的对应实验项目的“作业”模块内。
    \item “预习报告”还须拷贝到“实验报告”中（便以教师批改）。
    \item 大学物理实验”课程选做实验使用本“预习报告”；必做实验无须写预习报告，前往“学在浙大”完成预习测试即可。
\end{enumerate}

\vspace{1cm}
\xiaosi\centering \textbf{浙江大学物理实验教学中心制}

\end{document}